% Chapter 2, Figure 16: undamped response to a step disturbance in yaw (scan: figures/ch2/fig16.png).
% Curve: fig16.py, eq. (30) with eqs. (16), (17), (31a), (31b) and C_2 = 0: alpha_X = (M_s/C_1)(1 - cos omega_n t),
% in units of M_s/C_1 against omega_n t. First peak 2M_s/C_1 at pi/omega_n (eqs. (32a), (32b)); the dashed line is
% the offset axis M_s/C_1 it oscillates about (guide); the lines from the peak to the axes are solid as printed
% (peak line). One layout and time scale for Figs 16-19 (same axes, t to 2.75 pi/omega_n).
\documentclass[11pt]{standalone}
\usepackage{tamrfig}
\begin{document}
\begin{tikzpicture}
  \begin{axis}[tamr sketch,
      width=3.8in, height=2.5in,
      xmin=0, xmax=9.35, ymin=-0.36, ymax=2.62,
      xtick={3.14159}, xticklabels={$\dfrac{\pi}{\omega_n}$},
      ytick={1,2}, yticklabels={$\dfrac{M_s}{C_1}$, $\dfrac{2M_s}{C_1}$},
      xlabel={$t$ (sec)}, ylabel={$\alpha_X$ (rad)},
      % tamr sketch's rotate=-90 turns this label to read downward: reset it to upright, above the arrow
      every axis y label/.style={at={(ticklabel* cs:1.0)}, anchor=south}]
    \draw[guide] (axis cs:0,1) -- (axis cs:8.63938,1);
    \draw[peak line] (axis cs:0,2) -- (axis cs:3.14159,2) -- (axis cs:3.14159,0);
    \addplot[series1] table[x=t, y=alpha, col sep=comma] {figures/v2/ch2/fig16.csv};
    \node[anchor=east, font=\footnotesize, text=ink2] at (axis cs:0,0) {0};
  \end{axis}
\end{tikzpicture}
\end{document}
